%-------------------------
% Resume in Latex
% Author : Jake Gutierrez
% Based off of: https://github.com/sb2nov/resume
% License : MIT
%------------------------

\documentclass[letterpaper,10pt]{article}

\usepackage{latexsym}
\usepackage[empty]{fullpage}
\usepackage{titlesec}
\usepackage{marvosym}
\usepackage[usenames,dvipsnames]{color}
\usepackage{verbatim}
\usepackage{enumitem}
\usepackage[hidelinks]{hyperref}
\usepackage{fancyhdr}
\usepackage[english]{babel}
\usepackage{tabularx}
\usepackage{amssymb}
\usepackage{fontawesome5}
\usepackage{pifont}
\input{glyphtounicode}

\newcommand{\githubstars}[1]{[{\faGithubSquare}\,\ding{72}\,\textbf{#1}]}


%----------FONT OPTIONS----------
% sans-serif
% \usepackage[sfdefault]{FiraSans}
% \usepackage[sfdefault]{roboto}
% \usepackage[sfdefault]{noto-sans}
% \usepackage[default]{sourcesanspro}

% serif
% \usepackage{CormorantGaramond}
% \usepackage{charter}


\pagestyle{fancy}
\fancyhf{} % clear all header and footer fields
\fancyfoot{}
\renewcommand{\headrulewidth}{0pt}
\renewcommand{\footrulewidth}{0pt}

% Adjust margins
\addtolength{\oddsidemargin}{-0.55in}
\addtolength{\evensidemargin}{-0.55in}
\addtolength{\textwidth}{1.1in}
\addtolength{\topmargin}{-.65in}
\addtolength{\textheight}{1.3in}

\urlstyle{same}

\raggedbottom
\raggedright
\setlength{\tabcolsep}{0in}

% Sections formatting
\titleformat{\section}{
  \vspace{-4pt}\scshape\raggedright\large
}{}{0em}{}[\color{black}\titlerule \vspace{-4pt}]

% Ensure that generate pdf is machine readable/ATS parsable
\pdfgentounicode=1

%-------------------------
% Custom commands
\newcommand{\resumeItem}[1]{
  \item\small{
    {#1 \vspace{-5pt}}
  }
}

\newcommand{\resumeSubheading}[4]{
  \vspace{-3pt}\item
    \begin{tabular*}{0.97\textwidth}[t]{l@{\extracolsep{\fill}}r}
      \textbf{#1} & #2 \\
      \textit{\small#3} & \textit{\small #4} \\
    \end{tabular*}\vspace{-12pt}
}

\newcommand{\resumeSubSubheading}[2]{
    \item
    \begin{tabular*}{0.97\textwidth}{l@{\extracolsep{\fill}}r}
      \textit{\small#1} & \textit{\small #2} \\
    \end{tabular*}\vspace{-8pt}
}

\newcommand{\resumeProjectHeading}[2]{
    \item
    \begin{tabular*}{0.97\textwidth}{l@{\extracolsep{\fill}}r}
      \small#1 & #2 \\
    \end{tabular*}\vspace{-8pt}
}

\newcommand{\resumeSubItem}[1]{\resumeItem{#1}\vspace{-4pt}}

\renewcommand\labelitemii{$\vcenter{\hbox{\tiny$\bullet$}}$}

\newcommand{\resumeSubHeadingListStart}{\begin{itemize}[leftmargin=0.15in, label={}]}
\newcommand{\resumeSubHeadingListEnd}{\end{itemize}}
\newcommand{\resumeItemListStart}{\begin{itemize}}
\newcommand{\resumeItemListEnd}{\end{itemize}\vspace{-10pt}}

%-------------------------------------------
%%%%%%  RESUME STARTS HERE  %%%%%%%%%%%%%%%%%%%%%%%%%%%%


\begin{document}

%----------HEADING----------
\begin{center}
    \textbf{\LARGE \scshape Hyungtae Lim} \\ \vspace{1pt}
    \small
    \href{mailto:shapelim@mit.edu}{\underline{shapelim@mit.edu}} $|$
    \href{https://limhyungtae.github.io/hyungtae-lim/}{\underline{limhyungtae.github.io}} $|$
    \href{https://linkedin.com/in/hyungtae-lim}{\underline{linkedin.com/in/hyungtae-lim}} $|$
    \href{https://github.com/limhyungtae}{\underline{github.com/limhyungtae}}
\end{center}


% Skills and Experience before Education is recommended for industry applications.
% Reference: https://www.sweresume.app/articles/faang-resume-guide/

%-----------TECHNICAL SKILLS-----------
\section{Technical Skills}
 \begin{itemize}[leftmargin=0.15in, label={}]
    \small{\item{
     \textbf{Programming}{: C++, Python, CUDA, CMake} \\
     \textbf{Frameworks \& Tools}{: ROS/ROS2, PyTorch, OpenCV, Git, Docker} \\
     \textbf{Research Areas}{: SLAM, 3D LiDAR, Point Cloud Registration, Semantic Segmentation, Sensor Fusion, Deep Learning} \\
     \textbf{Languages}{: Korean (Native), English (Fluent)}
    }}
 \end{itemize}


%-----------EXPERIENCE-----------
\section{Experience}
  \resumeSubHeadingListStart

    \resumeSubheading
      {Postdoctoral Associate}{Apr. 2024 -- Present}
      {SPARK Lab, MIT (Advisor: Prof. Luca Carlone)}{Cambridge, MA}
      \resumeItemListStart
        \resumeItem{Reduced multi-session map inconsistency in city-scale environments by building a large-scale multi-robot SLAM system enabling real-time 3D mapping across ~8 robots and sessions.}
        \resumeItem{Accelerated research output by mentoring 5+ graduate and undergraduate researchers, contributing to 4+ published papers on 3D perception and autonomous mapping.}
        \resumeItem{Grew combined GitHub adoption to 4K+ stars by maintaining \href{https://github.com/MIT-SPARK/TEASER-plusplus}{TEASER++} \githubstars{2.1K}, \href{https://github.com/MIT-SPARK/Hydra}{Hydra} \githubstars{773}, \href{https://github.com/MIT-SPARK/KISS-Matcher}{KISS-Matcher} \githubstars{643}, \href{https://github.com/MIT-SPARK/VGGT-SLAM}{VGGT-SLAM} \githubstars{882}.}
      \resumeItemListEnd

    \resumeSubheading
      {Postdoctoral Fellow}{Mar. 2023 -- Mar. 2024}
      {Urban Robotics Lab, KAIST (Advisor: Prof. Hyun Myung)}{Daejeon, Korea}
      \resumeItemListStart
        \resumeItem{Improved localization robustness in dynamic scenes by developing real-time LiDAR-inertial and visual SLAM pipelines, deployed on 3+ mobile robot platforms.}
        \resumeItem{Reduced dynamic-object-induced map corruption by leveraging deep learning perception modules for moving object segmentation in crowded and unstructured environments.}
      \resumeItemListEnd

    \resumeSubheading
      {Visiting Scholar}{Nov. 2022 -- Feb. 2023}
      {StachnissLab, University of Bonn (Advisor: Prof. Cyrill Stachniss)}{Bonn, Germany}
      \resumeItemListStart
        \resumeItem{Cut dynamic-object artifacts in LiDAR maps by ${\sim}$90\% by building ERASOR2, an instance-aware removal system for static map building in urban scenes (RSS 2023).}
      \resumeItemListEnd

    \resumeSubheading
      {Research Intern}{Apr. 2022 -- Sep. 2022}
      {NAVER LABS}{Gyeonggi-Do, Korea}
      \resumeItemListStart
        \resumeItem{Achieved robust global LiDAR registration tolerating up to 99\% outlier correspondences by developing Quatro, exploiting ground segmentation for reliable loop closure in degenerate urban environments (ICRA 2022 / IJRR 2023).}
      \resumeItemListEnd

  \resumeSubHeadingListEnd


%-----------KEY OPEN-SOURCE PROJECTS-----------
% NOTE(hlim): These open-source projects with GitHub stars demonstrate coding quality
% and real-world adoption — leave this section fixed across applications.
% Stars signal that others rely on your code, which is exactly what production teams want.
\section{Key Open-Source Projects}
 \begin{itemize}[leftmargin=0.15in, label={}]
    \small{\item{
     \href{https://github.com/MIT-SPARK/KISS-Matcher}{\textbf{KISS-Matcher}} \githubstars{643}\emph{ICRA'25}  $|$ Out-of-the-box fast and robust global point cloud registration for large-scale mapping. \\
     \href{https://github.com/MIT-SPARK/VGGT-SLAM}{\textbf{VGGT-SLAM}} \githubstars{882} \emph{NeurIPS'25} $|$  -- Dense real-time RGB SLAM on SL(4) manifold; globally consistent 3D reconstruction. \\
     \href{https://github.com/url-kaist/patchwork-plusplus}{\textbf{Patchwork++}} \githubstars{842} \emph{IROS'22} $|$  $~$100\,Hz LiDAR ground segmentation; deployed in AV and mobile robot production. \\
     \href{https://github.com/LimHyungTae/ERASOR}{\textbf{ERASOR}} \githubstars{494} \emph{RA-L'21} $|$  -- ${\sim}$94\% dynamic-object artifact removal, enabling long-term autonomous navigation.
    }}
 \end{itemize}


%-----------KEY RELEVANT PROJECTS-----------
% NOTE(hlim): *** CUSTOMIZE THIS SECTION FOR EVERY APPLICATION. ***
% Unlike open-source above (fixed), swap these entries to match each job description.
%   Applying to AV perception team? -> Highlight 3D detection, segmentation, sensor fusion.
%   Applying to mapping/localization? -> Highlight SLAM, HD map, loop closure projects.
%   Applying to robot learning team?  -> Highlight RL, imitation learning, sim-to-real.
% XYZ formula: "Accomplished [X] as measured by [Y], by doing [Z]."
% Reference: https://www.sweresume.app/articles/faang-resume-guide/
\section{Key Relevant Projects}
    \resumeSubHeadingListStart
      \resumeProjectHeading
          {\href{https://github.com/your-org/project1}{\textbf{[e.g., Real-Time LiDAR SLAM System]}} $|$ \emph{[Category]}}{[Team / Year]}
          \resumeItemListStart
            \resumeItem{Reduced [X] by [Y]\% in city-scale environments by [Z]; deployed on [N]+ platforms at [scale].}
          \resumeItemListEnd
      \resumeProjectHeading
          {\href{https://github.com/your-org/project2}{\textbf{[e.g., 3D Object Detection Pipeline]}} $|$ \emph{[Category]}}{[Team / Year]}
          \resumeItemListStart
            \resumeItem{Improved [X] recall by [Y]\% by [Z], processing at [N]\,Hz on [hardware] (e.g., NVIDIA Orin).}
          \resumeItemListEnd
    \resumeSubHeadingListEnd


%-----------PUBLICATIONS-----------
\section{Publications}
 \begin{itemize}[leftmargin=0.15in, label={}]
    \small{\item{
     10\texttt{+} first-author papers at top-tier robotics and AI venues: IEEE ICRA, RSS, IEEE RA-L, NeurIPS, IEEE/CVF ICCV, IJRR
    }}
 \end{itemize}


%-----------HONORS \& AWARDS-----------
\section{Honors \& Awards}
 \begin{itemize}[leftmargin=0.15in, label={}]
    \small{\item{
     \textbf{2025} -- IEEE ICRA'25 Outstanding Reviewer Award (5 awardees out of 7,641 reviewers) \\
     \textbf{2024} -- RSS Pioneers 2024 (30 selected out of 202 candidates) \\
     \textbf{2024} -- 1st Prize, HILTI SLAM Challenge in IEEE ICRA'24 \\
     \textbf{2023} -- Best Paper Award, IEEE RA-L (5 of 1,100 papers) \\
     \textbf{2023} -- 1st Prize, HILTI SLAM Challenge in IEEE ICRA'23 (63 international teams) \\
     \textbf{2023} -- CES'23 Innovation Award (via technology transfer with HILLS Robotics)
    }}
 \end{itemize}


%-----------EDUCATION-----------
% (Moved below Skills & Experience per industry best practice:
%  https://www.sweresume.app/articles/faang-resume-guide/)
\section{Education}
  \resumeSubHeadingListStart
    \resumeSubheading
      {Korea Advanced Institute of Science and Technology (KAIST)}{Daejeon, Korea}
      {Ph.D. in Electrical Engineering and Robotics Program}{Mar. 2020 -- Feb. 2023}
    \resumeSubheading
      {Korea Advanced Institute of Science and Technology (KAIST)}{Daejeon, Korea}
      {M.S. in Electrical Engineering and Robotics Program}{Mar. 2018 -- Feb. 2020}
    \resumeSubheading
      {Korea Advanced Institute of Science and Technology (KAIST)}{Daejeon, Korea}
      {B.S. in Mechanical Engineering}{Mar. 2013 -- Feb. 2028}

  \resumeSubHeadingListEnd


%-------------------------------------------
\end{document}
